% TOP_PROBLEM: {"id": 5500034, "title": "Extending Pseudosegment Arrangements by Subdivision", "category_id": 2, "category": {"id": 2, "name": "combinatorics", "display_name": "Combinatorics", "description": "Counting problems, graph theory, discrete structures.", "slug": "combinatorics", "order_index": 2, "created_at": "2026-07-31T15:26:25.671Z"}, "status": "open"}
% FIRST_POSTED: null
% ENTRY_KIND: statement_only
% Imported from: batch07.tex; UnsolvedMath ID 5500034
\documentclass[11pt]{article}
\usepackage{amsmath,amssymb,mathtools}
\usepackage{fontspec,unicode-math}
\setmainfont{Latin Modern Roman}
\setsansfont{Latin Modern Sans}
\setmathfont{Latin Modern Math}
\usepackage{xurl,hyperref}
\newcommand{\sref}[2]{\hyperlink{source:#1:#2}{[S#2]}}
\newcommand{\cataloguescope}[1]{\paragraph{Mathematical question.}#1}
\begin{document}
% UnsolvedMath ID 5500034; source catalogue code AMR-054-0034
\setcounter{section}{5500033}
\section{Extending Pseudosegment Arrangements by Subdivision}\label{5500034}
{\small\sffamily UnsolvedMath ID 5500034 · Combinatorics}
\subsection{Definitions and mathematical statement}

\paragraph{Shared notation.}$\mathbb N=\{1,2,\ldots\}$, and $\mathbb Z,\mathbb Q,\mathbb R,\mathbb C$ denote the integers, rationals, reals and complex numbers. Any other conventions are specified locally.


A pseudosegment arrangement consists of $n$ finite planar curves with at most one intersection for each pair. A pseudoline arrangement consists of curves extending to infinity in both directions with at most one intersection for each pair. A subdivision turns selected intersection points into endpoints, splitting the incident pseudosegments. An extension continues the resulting pieces into pseudolines while preserving the arrangement of the pieces.
\paragraph{Statement recorded in the supplied source.}
Determine, as a function of $n$, the worst-case minimum number of intersections that must be selected as subdivision vertices to make a pseudosegment arrangement extendible to a pseudoline arrangement. \sref{5500034}{2}
\subsection{Short English statement}
Determine, as a function of $n$, the worst-case minimum number of intersections that must be selected as subdivision vertices to make a pseudosegment arrangement extendible to a pseudoline arrangement.
\subsection{Sources}
\begin{description}
\item[\hypertarget{source:5500034:1}{S1}] ULAM.AI and the UnsolvedMath Contributors, UnsolvedMath: A Curated Collection of Open Mathematics Problems, record 5500034 (AMR-054-0034). CC BY 4.0. \url{https://huggingface.co/datasets/ulamai/UnsolvedMath}.
\item[\hypertarget{source:5500034:2}{S2}] Erik D. Demaine, Joseph S. B. Mitchell and Joseph O\textquotesingle Rourke (editors), \emph{The Open Problems Project}, maintained original source, Problem 34, Statement and complete Partial/Related Results. \url{https://github.com/edemaine/topp/blob/main/Problems/P.000034}.
\end{description}
\label{end:5500034}
\end{document}
